% Part: lambda-calculus
% Chapter: introduction
% Section: beta

\documentclass[../../../include/open-logic-section]{subfiles}

\begin{document}

\olfileid{lam}{int}{bet}
\olsection{$\beta$-تقلیل}

جدوں $(\lambd[m][(\lambd[y][y]) m])$ ویکھیے، تاں سُبھاوک خیال
آؤندا اے کہ ایہدا $\lambd[m][m]$ نال کوئی جوڑ اے، یعنی پہلا
ٹرم «سوکھا» کرن نال دوجا ٹرم ملنا چاہیدا اے۔ \emph{$\beta$-تقلیل}
ایس خیال نوں باقاعدہ شکل دیندی اے۔

\begin{defn}[$\beta$-سکڑاؤ، $\bredone$] \ollabel{defn:betacontr}
  \emph{$\beta$-سکڑاؤ} ($\bredone$) ٹرماں اُتے سبھ توں چھوٹا
  موافق تعلق اے جیہڑا ہِٹھلی شرط پوری کردا اے۔
  \[
    (\lambd[x][N])Q \bredone \Subst{N}{Q}{x}
  \]
  اسیں آکھدے آں کہ $P$، \emph{$\beta$-سکڑاؤ} نال $Q$ بن جاندا
  اے جے $P \bredone Q$۔ $(\lambd[x][N])Q$ دی شکل والے ٹرم
  نوں \emph{ریڈیکس} آکھدے آں۔
\end{defn}

\begin{prob} \ollabel{prob:def}
  $\beta$-سکڑاؤ دی ہم ارز استقرائی تعریف کھول کے لکھو، جیویں
  اسیں پابند متغیر دا ناں بدلن لئی \olref[lam][syn][alp]{defn:aconvone}
  وچ کیتا سی۔
\end{prob}

\begin{defn}[$\beta$-تقلیل، $\bred$] \ollabel{defn:betared}
  \emph{$\beta$-تقلیل} ($\bred$) ٹرماں اُتے سبھ توں چھوٹا
  انعکاسی تے تعدّی والا تعلق اے جیہدے وچ $\bredone$ شامل اے۔
  اسیں آکھدے آں کہ $P$ دی $\beta$-تقلیل نال $Q$ ملدا اے جے
  $P \bred Q$۔
\end{defn}

جدوں گل دا مطلب صاف ہووے، اسیں $\redone$ لکھاں گے
$\bredone$ دی تھاں، تے $\red$ لکھاں گے $\bred$ دی تھاں۔

سادے لفظاں وچ، $M \bred N$ اُدوں تے صرف اُدوں اے جدوں $M$
نوں $N$ وچ $\beta$-سکڑاؤ دے صفر یا کئی قدماں نال بدلیا جا سکے۔

\begin{defn}[$\beta$-نارمل]
جیہڑا ٹرم $\beta$-سکڑاؤ نال ہور نہیں سکڑ سکدا، اوہنوں
\emph{$\beta$-نارمل} آکھدے آں۔
\end{defn}

جے $M \bred N$ تے $N$، $\beta$-نارمل اے، تاں اسیں
$N$ نوں~$M$ دا \emph{نارمل روپ} آکھدے آں۔ کوئی پُچھ سکدا اے
کہ کی ٹرم دا نارمل روپ اکّو ہوندا اے۔ جواب ہاں اے، جیویں
اگّے ویکھاں گے۔

کجھ مثالاں ویکھیے۔
\begin{enumerate}
\item ساڈے کول ایہ تقلیلاں نیں۔
\begin{align*}
(\lambd[x][xxy]) \lambd[z][z] & \redone (\lambd[z][z])(\lambd[z][z]) y \\
& \redone (\lambd[z][z]) y \\
& \redone y
\end{align*}
\item ٹرم نوں «سوکھا» کرنا اوہنوں ہور گُنجھلدار وی بنا سکدا اے۔
\begin{align*}
(\lambd[x][xxy])(\lambd[x][xxy]) & \redone (\lambd[x][xxy])(\lambd[x][xxy])y \\
& \redone (\lambd[x][xxy])(\lambd[x][xxy])yy \\
& \redone \dots
\end{align*}
\item ایہ ٹرم نوں اویویں وی رہن دے سکدا اے۔
\[
(\lambd[x][xx])(\lambd[x][xx]) \redone (\lambd[x][xx])(\lambd[x][xx])
\]
\item کجھ ٹرماں دی اک توں ودھ طریقیاں نال تقلیل ہو سکدی اے۔
  مثال لئی، سبھ توں باہرلے اطلاق دے سکڑاؤ نال ایہ ملدا اے۔
\[
(\lambd[x][(\lambd[y][yx]) z]) v \redone (\lambd[y][yv]) z
\]
تے سبھ توں اندرلے اطلاق دے سکڑاؤ نال ایہ ملدا اے۔
\[
(\lambd[x][(\lambd[y][yx]) z]) v \redone (\lambd[x][zx]) v
\]
دھیان دیو کہ ایس صورت وچ دوہاں ٹرماں دی اگلی تقلیل نال
اکّو ٹرم $zv$ ملدا اے۔
\end{enumerate}

پچھلی مثال دا آخری نتیجہ اتفاق نہیں۔ ایہ لامبڈا کیلکولس دی
اک ڈونگھی تے اہم خاصیت دی مثال اے جیہنوں چرچ–روسر خاصیت
آکھدے آں۔

\begin{digress}
  عام طور اُتے ٹرم دی $\beta$-تقلیل دا اک توں ودھ طریقہ ہوندا
  اے۔ ایس لئی تقلیل دیاں کئی چالاں بنائیاں گئیاں نیں، جنہاں
  وچوں سبھ توں عام \emph{سُبھاوک چال} اے۔ سُبھاوک چال ہمیشہ
  \emph{سبھ توں کھبے} ریڈیکس نوں سکیڑدی اے۔ ریڈیکس دی تھاں
  ٹرم وچ اوہدی شروعات والی تھاں اے۔ ایس چال دی خاصیت ایہ اے
  کہ کسے چال نال ٹرم دی تقلیل نارمل روپ تک اُدوں تے صرف اُدوں
  ہو سکدی اے جدوں سُبھاوک چال نال اوہدی تقلیل نارمل روپ تک
  ہو سکے۔ اگّے اسیں سُبھاوک چال ورتاں گے، جے کوئی ہور چال نہ
  دسی جاوے۔
\end{digress}

\begin{defn}[$\beta$-ہم ارزی، $\equal$]
  \emph{$\beta$-ہم ارزی} ($\equal$) اوہ تعلق اے جیہڑا استقرا نال
  ایویں متعین ہوندا اے۔
  \begin{enumerate}
  \item $M \equal M$۔
  \item جے $M \equal N$، تاں $N \equal M$۔
  \item جے $M \equal N$ تے $N \equal O$، تاں $M \equal O$۔
  \item جے $M \equal N$، تاں $PM \equal PN$۔
  \item جے $M \equal N$، تاں $MQ \equal NQ$۔
  \item جے $M \equal N$، تاں $\lambd[x][M] \equal \lambd[x][N]$۔
  \item $(\lambd[x][N])Q \equal \Subst{N}{Q}{x}$۔
  \end{enumerate}
\end{defn}

پہلے تِن قاعدے ایس تعلق نوں ہم ارزیتی تعلق بناؤندے نیں۔
اگلے تِن ایہنوں موافق بناؤندے نیں۔ آخری قاعدہ یقین بناؤندا اے
کہ $\beta$-سکڑاؤ ایس وچ شامل اے۔

سادے لفظاں وچ، دو ٹرم $\beta$-ہم ارز اُدوں تے صرف اُدوں
ہوندے نیں جدوں اک نوں دوجے وچ $\beta$-سکڑاؤ یا
$\beta$-سکڑاؤ دے «اُلٹ» دے صفر یا ودھ قدماں نال بدلیا جا
سکے۔ $\beta$-سکڑاؤ دا اُلٹ ایویں متعین اے کہ $M$ دا
اُلٹا-$\beta$-سکڑاؤ $N$ تک اُدوں تے صرف اُدوں اے جدوں $N$
دا $\beta$-سکڑاؤ $M$ تک اے۔

اُپرلے قاعدیاں دے نال اسیں ہور قاعدے پا کے تعلق نوں ودھاواں گے
تے ودھے ہوئے ہم ارزیتی تعلق نوں $\equal[X]$ لکھاں گے، جتھے $X$
وادھا کرن والا قاعدہ اے۔

\end{document}
